\capitulo{8}{Lineas de trabajo futuras}
El presente proyecto ha logrado satisfacer con creces los objetivos planteados, culminando en una herramienta de grado de producción completamente funcional para la unidad (\textit{URCCPQ}). No obstante, el diseño de la arquitectura de \textit{software} se ha concebido bajo el principio de escalabilidad, sentando unas bases sólidas que permiten futuras iteraciones hacia un ecosistema de integración clínica absoluta.

En cuanto al flujo de información, la línea de trabajo futuro más ambiciosa y con mayor impacto clínico es la automatización total de la ingesta de datos. En su versión actual, la herramienta requiere que los facultativos introduzcan manualmente los registros clínicos en la base de datos \textit{PostgreSQL} a través de los formularios de la interfaz. El objetivo a medio plazo es desarrollar conectores directos a nivel de base de datos o módulos de interoperabilidad (\textit{APIs}) que integren la aplicación directamente con el sistema de información clínica \textit{ICCA} (\textit{IntelliSpace Critical Care and Anesthesia}) utilizado en la unidad. Esta integración permitiría alimentar el motor algorítmico de \textit{Python} en tiempo real extrayendo la información de origen, calculando los indicadores de la \textit{SEMICYUC} de forma continua y eliminando por completo la necesidad de transcripción o intervención manual por parte del personal médico.

Finalmente, la consolidación de esta conexión automatizada con \textit{ICCA} abriría la puerta a expandir el motor estadístico actual. Aprovechando librerías como \textit{Scikit-Learn}, ya integradas en el proyecto, se podrían entrenar modelos predictivos de \textit{Machine Learning}\footnote{Rama de la Inteligencia Artificial que permite a las máquinas aprender y mejorar de forma autónoma a través de datos y algoritmos, sin ser programadas explícitamente.} capaces de alertar sobre desviaciones en los indicadores antes de que ocurran, transformando la aplicación de una herramienta de monitorización y análisis a un verdadero sistema inteligente de soporte a la decisión clínica.

En el Anexo C.4: Instrucciones para la modificación o mejora del proyecto,
se incluyen comentarios de interés y se detallan a fondo cada uno de los
aspectos susceptibles de mejora mencionados en este apartado.